% helix_wave_timing_derivation.tex
% Self-contained derivation for the HELIX plate-impact timing model.
\documentclass[11pt]{article}

\usepackage[T1]{fontenc}
\usepackage[utf8]{inputenc}
\usepackage[margin=1in]{geometry}
\usepackage{amsmath,amssymb,mathtools}
\usepackage{hyperref}
\usepackage{xparse}

\hypersetup{
  colorlinks=true,
  linkcolor=blue,
  urlcolor=blue,
  pdftitle={One-dimensional plate-impact wave timing derivation}
}

\newcommand{\rhoz}[1][]{\rho_{0#1}}
\newcommand{\Us}[1][]{U_{\mathrm{s}#1}}
\newcommand{\up}[1][]{u_{\mathrm{p}#1}}
\newcommand{\sig}{\sigma}
\NewDocumentCommand{\CL}{o}{C^{\mathrm L}\IfValueT{#1}{_{#1}}}
\NewDocumentCommand{\CLv}{om}{C^{\mathrm L,(#2)}\IfValueT{#1}{_{#1}}}
\newcommand{\tFR}{t_{\mathrm{FR}}^{(0)}}
\newcommand{\Pt}{P_t}

\title{One-Dimensional Plate-Impact Kinematics and Wave-Timing Derivation}
\author{For the HELIX/SPADE diagnostic timing overlay}
\date{\today}

\begin{document}
\maketitle

\begin{abstract}
This note derives the equations needed to infer plate-impact kinematics from a
measured free-surface velocity and to identify what timing information can,
and cannot, be obtained from a constant-speed construction.  It distinguishes
(i) results that follow directly from the Rankine--Hugoniot jump conditions
from (ii) release and reflection predictions that require a release model.
The flyer-release free-flight arrival and the target free-surface reverberation
are shown to have mutually exclusive isolated-path conditions.  Consequently,
neither a universal no-spall plateau duration nor a universal recompression
series follows from only $C_0$, $S$, two thicknesses, and one nominal sound
speed.  A finite-amplitude characteristic calculation is stated as the
quantitative reference, followed by explicitly labelled first-hand timing
markers for incomplete target data.  The jump-condition and release concepts
used here follow the standard
shock-compression treatment, while the stated limitations follow directly from
release-adiabat and release-wave measurements \cite{Rice1958,Kerley2008,Hawreliak2024}.
\end{abstract}

\tableofcontents

\section{Scope, geometry, and notation}

Consider planar normal impact of a flyer (subscript $f$) on a target
(subscript $t$).  The target is initially at rest, the flyer initially moves
in the positive $x$ direction at speed $V$, and the flyer--target interface is
at $x=0$ at first contact.  The target free surface is at $x=h_{t0}$ and the
flyer rear free surface is at $x=-h_{f0}$.  The symbols $h_{f0}$ and $h_{t0}$
are the \emph{initial} (unshocked) thicknesses.

Compression and longitudinal normal stress are taken positive.  The
derivation assumes one-dimensional, planar waves, perfect contact, zero
initial stress, and no spall in the base solution.  A real elastic--plastic
solid can additionally have an elastic precursor, strength effects, and a
nonlinear release; those are discussed as limitations below
\cite{Rice1958,Hawreliak2024}.

For material $i\in\{f,t\}$, let
\begin{equation}
  \rhoz[i],\qquad C_{0i},\qquad S_i
\end{equation}
denote its initial density, linear-Hugoniot intercept, and linear-Hugoniot
slope.  The assumed principal Hugoniot is
\begin{equation}
  \Us[i](\up[i])=C_{0i}+S_i\up[i].
  \label{eq:linear-hugoniot}
\end{equation}
This $\Us$--$\up$ form is an empirical representation of the principal
Hugoniot, used here only over the range in which its material parameters have
been validated \cite{Kerley2006}.
Here $\Us[i]$ is the shock speed \emph{relative to the material ahead of the
shock}, and $\up[i]$ is the particle-velocity jump across that shock.  Thus,
for the target shock, the laboratory shock speed is $\Us[t]$; for the shock
moving back through the moving flyer, its laboratory speed is $V-\Us[f]$.

The jump-condition stress in this scalar model is
\begin{equation}
  \sig_H(\up)=\rhoz\Us(\up)\up
  =\rhoz\bigl(C_0+S\up\bigr)\up.
  \label{eq:hugoniot-stress}
\end{equation}
For a solid this is more precisely the longitudinal Hugoniot stress in a
one-dimensional strength-free approximation; calling it a hydrostatic
``pressure'' is conventional shorthand, not an additional constitutive law
\cite{Rice1958}.

\section{Interface state and inverse impact-velocity reconstruction}

\subsection{Impedance matching}

After the initial shocks have formed, both sides of the impact interface have
the same particle velocity $u$ and the same normal stress.  The target has
been accelerated from $0$ to $u$, whereas the flyer has been decelerated from
$V$ to $u$.  Therefore the two shock particle-velocity jumps are
\begin{equation}
  \up[t]=u,\qquad \up[f]=x=V-u.
  \label{eq:jumps}
\end{equation}
Stress continuity gives the exact impedance-match equation within the linear
Hugoniot model,
\begin{align}
  \rhoz[t]\bigl(C_{0t}+S_tu\bigr)u
  &=\rhoz[f]\bigl(C_{0f}+S_fx\bigr)x,
  \label{eq:impedance-match}\\
  V&=u+x.
\end{align}
This equation is the appropriate relation whether the flyer and target are
the same material or different materials \cite{Rice1958,Kerley2008}.

\subsection{What a free-surface velocity measures}

Let $U_{fs}$ be the free-surface particle velocity following release of the
target shock state.  In general, $U_{fs}$ is \emph{not} exactly $2u$.  For an
isentropic release from the shocked state to zero longitudinal stress,
\begin{equation}
  U_{fs}(u)=u+\Delta u_{\rm rel}(u),\qquad
  \Delta u_{\rm rel}(u)=
  \int_{0}^{\sig_H(u)}\frac{d\sig}{\rho(\sig)\,c(\sig)}.
  \label{eq:free-surface-exact}
\end{equation}
The local (spatial/Eulerian) longitudinal sound speed on the release is $c$,
and the integral requires a release adiabat (and, for a strength model, the
relevant longitudinal incremental modulus).  A Hugoniot alone does not supply
this information \cite{LyzengaAhrens1978,Kerley2008}.

The common first-order approximation is
\begin{equation}
  \Delta u_{\rm rel}\simeq u,
  \qquad\text{hence}\qquad u\simeq\frac{U_{fs}}{2}.
  \label{eq:free-surface-double}
\end{equation}
It is appropriate only when the measured peak is an unperturbed free-surface
state and the release is sufficiently close to the loading relation.  It is
not an exact finite-shock identity.  For some common metal states the exact
release gives $U_{fs}>2u$, in which case the doubling approximation
underestimates both $u$ and the inferred $V$; this sign is not universal and
must not be assumed without a material-specific release model
\cite{Kerley2008}.

\subsection{Closed-form inverse under the doubling approximation}

Given a measured $U_{fs}$, set $u=U_{fs}/2$ under
Eq.~\eqref{eq:free-surface-double} and calculate
\begin{equation}
  \sig_t=\rhoz[t]\bigl(C_{0t}+S_tu\bigr)u.
  \label{eq:target-stress}
\end{equation}
The flyer jump $x$ follows from
\begin{equation}
  \rhoz[f]S_f x^2+\rhoz[f]C_{0f}x-\sig_t=0.
\end{equation}
The physical root is
\begin{equation}
  x=
  \frac{-\rhoz[f]C_{0f}+
  \sqrt{(\rhoz[f]C_{0f})^2+4\rhoz[f]S_f\sig_t}}
  {2\rhoz[f]S_f},
  \qquad V=u+x.
  \label{eq:inverse-velocity}
\end{equation}
For $\rhoz[f]>0$, $C_{0f}>0$, and $S_f>0$, the left-hand side of the
quadratic is strictly increasing for $x\geq0$, since
\begin{equation}
  \frac{d}{dx}\left[\rhoz[f](C_{0f}+S_fx)x\right]
  =\rhoz[f](C_{0f}+2S_fx)>0.
\end{equation}
Thus a compressive target stress has exactly one non-negative flyer-jump root;
the other root is negative and does not represent a compressive shock in the
flyer.  Equations \eqref{eq:target-stress}--\eqref{eq:inverse-velocity} are
algebraically exact once Eq.~\eqref{eq:free-surface-double} and the two linear
Hugoniots have been adopted.

\section{Compression, sound speed, and the coordinate used for timing}
\label{sec:speed-coordinates}

The Rankine--Hugoniot mass jump condition gives the density immediately behind
a shock:
\begin{equation}
  \rho_H=\rhoz\frac{\Us}{\Us-\up}.
  \label{eq:shocked-density}
\end{equation}
Physical compression requires $\Us-\up>0$.  This condition must be checked
whenever arbitrary Hugoniot parameters are accepted as input.  For the linear
form,
\begin{equation}
  \Us-\up=C_0+(S-1)\up.
  \label{eq:compression-check-linear}
\end{equation}
If $S<1$, this expression becomes non-positive at
$\up=C_0/(1-S)$; hence $S\geq1$ is necessary for this particular linear form
to maintain $\Us>\up$ for every $\up\geq0$ (assuming $C_0>0$).  This does not
establish that a material follows a linear Hugoniot, and a fitted linear form
must remain restricted to its validated range \cite{Kerley2006}.

It is essential not to mix Eulerian and Lagrangian wave speeds.  Let $c$ be
the local/comoving sound speed measured relative to the shocked material (the
spatial/Eulerian sound speed), or a selected constant approximation to a
release-characteristic speed.  A characteristic in that state has laboratory
speed $u\pm c$.  Written in the unshocked material coordinate $X$, its
magnitude is instead
\begin{equation}
  \CL=\frac{\rho_H}{\rhoz}c
      =\frac{\Us}{\Us-\up}c.
  \label{eq:lagrangian-speed}
\end{equation}
Consequently, a traversal of an initial thickness $h_0$ takes
\begin{equation}
  t=\frac{h_0}{\CL}
   =\frac{h_0(\Us-\up)}{\Us\,c},
  \label{eq:transit-time}
\end{equation}
not generally $h_0/c$.

For a true finite-amplitude release, $c$ changes along the release fan and the
transit time must be calculated by the method of characteristics (MOC).  In a
first-order diagnostic overlay, choose one representative $c$ or $\CL$ and
state explicitly which is stored in the material-property table
\cite{CourantFriedrichs1948,Hawreliak2024}.

\section{Finite-amplitude characteristic calculation: the quantitative reference}
\label{sec:quantitative-reference}

The physically complete prediction is not a single round-trip formula.  Let
$H_i$ denote the shocked state of material $i$, and let
$\sig_{\mathrm{rel},i}(u;H_i)$ be its release adiabat from that state.  The
Lagrangian characteristic speed on that release is
\begin{equation}
  C_{\mathrm{rel},i}^{\mathrm L}(u;H_i)
  =\frac{\rho_i(u)}{\rhoz[i]}c_{\mathrm{inc},i}(u)
  =\frac{1}{\rhoz[i]}
    \left|\frac{d\sig_{\mathrm{rel},i}}{du}\right|_S .
  \label{eq:release-characteristic-speed}
\end{equation}
Here $c_{\mathrm{inc},i}$ is the appropriate local incremental speed: the
hydrodynamic sound speed for a fluid model, or the longitudinal tangent speed
when strength is retained.  The absolute value is required because the sign of
$d\sig/du$ depends on the characteristic family.  The material-coordinate
characteristics satisfy
\begin{equation}
  \frac{dX}{dT}=\pm C_{\mathrm{rel},i}^{\mathrm L}(u;H_i).
  \label{eq:release-characteristics}
\end{equation}
Consequently, the changing characteristic speed across a release fan, and the
state produced when it meets another wave, must be propagated by the method of
characteristics (MOC) or an equivalent finite-amplitude wave solver
\cite{CourantFriedrichs1948,Davison2008}.

In that calculation the target free surface imposes $\sig=0$.  At a closed
flyer--target interface, stress and particle velocity are continuous; for an
unbonded interface, the contact condition is unilateral and cannot sustain
tension.  The initial target free-surface release, the flyer rear-surface
release, their interaction, and any interface reflection are therefore solved
together.  If $\mathcal R_{\mathrm{rc}}$ is the family of target
characteristics that reaches the target free surface as a returned compression,
the first such arrival is the solver output
\begin{equation}
  t_{\mathrm{rc}}=\min_{\chi\in\mathcal R_{\mathrm{rc}}}
  \left[T_\chi(X=h_{t0})-T_s\right].
  \label{eq:exact-recompression-arrival}
\end{equation}
Likewise, the first departure from the initial free-surface state is the
earliest arrival among \emph{all} wave families at that boundary.  These
definitions replace a universal ``plateau duration'' or an $n\Pt$ series;
they reduce to the constant-speed isolated-path markers only under the
restrictions stated later.

\subsection{Data required for the quantitative estimate}

The MOC calculation requires more than principal-Hugoniot coefficients.  At a
minimum it needs:
\begin{itemize}
  \item the initial densities, thicknesses, impact speed, and initial states;
  \item an EOS or complete Hugoniot description for \emph{both} flyer and
        target, sufficient to determine the initial interface state;
  \item release adiabats or equivalent state-dependent incremental moduli for
        both materials, so that Eq.~\eqref{eq:release-characteristic-speed}
        is known along each release path;
  \item an elastic--plastic/strength model whenever an elastic precursor or
        longitudinal release head is material to the measurement;
  \item free-surface and interface boundary conditions, including a unilateral
        contact/separation rule for an unbonded flyer; and
  \item a distinct damage/spall model if a later pullback or recompaction is
        to be interpreted as a spall signal.
\end{itemize}
A linear $\Us$--$\up$ relation supplies only the principal-Hugoniot loading
locus; it does not supply the release adiabat in
Eq.~\eqref{eq:release-characteristic-speed}
\cite{LyzengaAhrens1978,Kerley2008}.  Thus the output of this section is the
reference calculation against which a timing overlay should be judged.

\section{Flyer-release free-flight arrival and shock--release catch-up}

Let global time $T=0$ be first flyer--target contact.  The target shock reaches
the target free surface at
\begin{equation}
  T_s=\frac{h_{t0}}{\Us[t]}.
  \label{eq:shock-arrival-fs}
\end{equation}
The flyer shock reaches the flyer rear surface after $h_{f0}/\Us[f]$.  The
rear-surface release then traverses the shocked flyer and reaches the impact
interface at
\begin{equation}
  T_i=\frac{h_{f0}}{\Us[f]}+\frac{h_{f0}}{\CL[f]}.
  \label{eq:flyer-release-interface}
\end{equation}
After entering the already shocked target, that release characteristic reaches
the target free surface at
\begin{equation}
  T_r=T_i+\frac{h_{t0}}{\CL[t]}.
  \label{eq:release-arrival-fs}
\end{equation}
If this flyer-generated release characteristic is allowed to travel without
interacting with the release launched at the target free surface, its arrival
time relative to target shock arrival is
\begin{align}
  \tFR
  &=T_r-T_s \notag\\
  &=\frac{h_{f0}}{\Us[f]}+\frac{h_{f0}}{\CL[f]}
    +h_{t0}\left(\frac{1}{\CL[t]}-\frac{1}{\Us[t]}\right).
  \label{eq:flyer-release-free-flight}
\end{align}
The superscript $(0)$ emphasizes that this is a \emph{free-flight} ray result,
not in general an observed plateau duration.  The first two terms are the
flyer round-trip contribution; the final term compares the target transit of
that release to the earlier target transit of the shock.  It may be negative or
positive, and is negative only when $\CL[t]>\Us[t]$.  Choosing two candidate
target speeds $\CLv[t]{1}$ and $\CLv[t]{2}$ changes this contribution by
\begin{equation}
  h_{t0}\left(\frac{1}{\CLv[t]{2}}-
                \frac{1}{\CLv[t]{1}}\right).
  \label{eq:target-speed-timing-spread}
\end{equation}
Thus even its sign cannot be inferred reliably by substituting an arbitrary
ambient bulk or longitudinal tabulation for the shocked-state release-head
speed.  In particular, comparing the shocked-state bulk and longitudinal
choices gives
\begin{equation}
  \Delta t_t^{(L-B)}=
  h_{t0}\frac{\rhoz[t]}{\rho_{H,t}}
  \left(\frac{1}{c_L}-\frac{1}{c_B}\right),
  \label{eq:longitudinal-bulk-timing-spread}
\end{equation}
where each speed is evaluated at the selected shocked state.  This sensitivity
is distinct from the finite-width release-fan uncertainty.  This is an isolated-ray,
constant-speed construction.  It should not be taken as validation of the
release-wave-overtake method: shocked-copper measurements show that the
front-surface-impact and release-wave-overtake methods can diverge when the
release response is complex \cite{Hawreliak2024}.

If Eq.~\eqref{eq:flyer-release-free-flight} is non-positive, a release characteristic
overtakes the primary shock before the target free surface.  The presumed
constant shocked state never produces a sustained free-surface plateau, so a
plateau/recompression overlay based on that state is not physically valid.
The catch-up point, when $\CL[t]>\Us[t]$, is
\begin{equation}
  X_c=\frac{\Us[t]\CL[t]T_i}{\CL[t]-\Us[t]}.
  \label{eq:catchup}
\end{equation}
The release catches the shock inside the target if $X_c<h_{t0}$, which is
equivalent to $\tFR<0$.

\section{Ordering of the two release paths: no universal plateau marker}
\label{sec:release-ordering}

The initial target shock reaching the free surface immediately launches a
left-going release.  Its isolated target round-trip time is
\begin{equation}
  \Pt=\frac{2h_{t0}}{\CL[t]},
  \label{eq:target-round-trip}
\end{equation}
and it would reach the impact interface at the global time
\begin{equation}
  T_{\mathrm{FS}\to I}=T_s+\frac{h_{t0}}{\CL[t]}.
  \label{eq:fs-release-interface-time}
\end{equation}
Combining Eqs.~\eqref{eq:flyer-release-free-flight},
\eqref{eq:target-round-trip}, and \eqref{eq:fs-release-interface-time} gives
the identity
\begin{equation}
  \tFR-\Pt=T_i-T_{\mathrm{FS}\to I}.
  \label{eq:ordering-identity}
\end{equation}
This identity exposes the structural limitation of a simple timing overlay.
Finite-amplitude interactions must be propagated with a characteristic or
equivalent wave-interaction calculation \cite{CourantFriedrichs1948,Davison2008}.

\paragraph{Case A: $\tFR<\Pt$.}
By Eq.~\eqref{eq:ordering-identity}, $T_i<T_{\mathrm{FS}\to I}$: the flyer
rear-surface release has entered the target before the initial free-surface
release reaches the interface.  The counter-propagating release fans meet in
the target.  In the strictly linear, constant-speed construction, those waves
superpose and $\tFR$ remains the free-flight arrival time.  At finite amplitude
their interaction needs a characteristic calculation; independently, the
flyer release has changed the interface state before the initial free-surface
release arrives there, so the unperturbed-state reflection coefficient cannot
generate an $n\Pt$ series.

\paragraph{Case B: $\tFR>\Pt$.}
The initial free-surface release reaches an interface that is still in the
original shocked state.  Its isolated reflection can return to the free surface
at $\Pt$, before the flyer release arrives.  Thus even in this special branch,
the free surface cannot be called a flat plateau through $\tFR$.

\paragraph{Equality.}
At $\tFR=\Pt$, the relevant wave interactions coincide.  It is a
multi-characteristic interaction.  Within the strictly linear, constant-speed
construction, the first free-surface departure time is exactly
\begin{equation}
  t_{\mathrm{first}}^{(0)}=\min\bigl(\tFR,\Pt\bigr).
  \label{eq:first-free-surface-departure}
\end{equation}
For finite-amplitude releases this equation is a diagnostic estimate, not a
substitute for a characteristic solution.

\section{Acoustic reflection signs in the isolated-path limit}

\subsection{Impedance appropriate to a small release/recompression}

The reflection of a small perturbation about a shocked state is governed by
the incremental acoustic impedance, not the Rayleigh (shock) impedance
$\rhoz\Us$.  With the notation of Section~\ref{sec:speed-coordinates},
\begin{equation}
  Z=\rho_Hc=\rhoz\CL.
  \label{eq:acoustic-impedance}
\end{equation}
Both forms in Eq.~\eqref{eq:acoustic-impedance} are identical only when
$c$ and $\CL$ are related by Eq.~\eqref{eq:lagrangian-speed}.

Suppose a small stress perturbation travels from the shocked target into the
shocked flyer.  Continuity of stress and particle velocity at the interface
gives the stress-amplitude reflection coefficient
\begin{equation}
  R_\sig=\frac{Z_f-Z_t}{Z_f+Z_t},
  \qquad
  \delta\sig_{\rm reflected}=R_\sig\,\delta\sig_{\rm incident}.
  \label{eq:reflection-coefficient}
\end{equation}
This is the standard one-dimensional, small-amplitude stress reflection
coefficient for incidence from the target side \cite{Davison2008}.
The first release created at a free surface has a stress change of order
$\sig_H$, not necessarily a small perturbation.  Applying
Eq.~\eqref{eq:reflection-coefficient} to that wave therefore supplies only a
linearized sign/topology guide; its finite-amplitude reflected state requires
the release EOS and characteristic construction.
If $Z_f<Z_t$, then $R_\sig<0$: a target-side incident release
($\delta\sig<0$) reflects as a compression, whereas a target-side incident
compression reflects as a release.

\subsection{Initial-free-surface release: the isolated recompression path}

At $t=0$ (where $t$ is now referenced to target shock arrival at the free
surface), the initial target compression reaches the free boundary.  It
reflects there as a release, which travels back to the impact interface.  For
$Z_f<Z_t$, Eq.~\eqref{eq:reflection-coefficient} turns that left-going release
into a right-going compression.  It returns to the target free surface after
one target round trip:
\begin{equation}
  t_{\rm rc,n}=n\Pt,\quad n=1,2,\ldots
  \label{eq:initial-fs-reverberation}
\end{equation}
In the constant-speed, small-perturbation and isolated-interface limit, these
are low-impedance-flyer \emph{recompression} markers referenced to the initial
shock arrival.  They are not a default series for a general flyer/target
experiment; Case~A of Section~\ref{sec:release-ordering} invalidates this
construction.

\subsection{Flyer-release return: a distinct isolated unloading path}

In free flight, the flyer rear-surface release reaches the target free surface
at $t=\tFR$.  At that free boundary it reflects as a compression that
travels back toward the impact interface.  If $Z_f<Z_t$, that compression
reflects from the interface as a release.  Its next target-free-surface
arrival is therefore
\begin{equation}
  t=\tFR+\Pt,
  \label{eq:flyer-release-return}
\end{equation}
and, in this linear sign analysis, it is an \emph{unloading} arrival rather
than a recompression.  Thus $\tFR+n\Pt$ must not be plotted as a
low-impedance recompression merely because it occurs after a pullback.

\paragraph{Contact condition.} An unbonded flyer--target interface is
unilateral: it may carry compression but not tension, $\sig_{xx}\geq0$ at
contact.  The isolated first reflection can leave a compressive interface
stress, but this must be checked after every finite-amplitude interaction.  A
``perfect contact'' boundary condition must not silently permit tensile
traction.

\section{Implementation-ready constant-speed model}

For a transparent first-order implementation, store or estimate the
shocked-state \emph{release-head} speed $c_{\mathrm{head},i}$ in the local
comoving frame.  The onset of unloading is governed by that leading
characteristic.  In an elastic--plastic solid it is generally a longitudinal
release-head speed; later portions of the release fan can propagate at a
different speed.  For an isotropic solid,
\begin{equation}
  c_L^2=c_B^2+\frac{4}{3}c_S^2,
  \label{eq:longitudinal-bulk-shear}
\end{equation}
so a bulk-wave value and a longitudinal-wave value should not be used
interchangeably.  In a fluid, $c_L=c_B$.  The relevant speeds must be those of
the \emph{shocked state}, not ambient tabulations.  Measured copper release
experiments explicitly use the leading release arrival to determine the
shocked-state Lagrangian longitudinal speed.  Hawreliak \emph{et al.} also
report a quasielastic release in copper at 190~GPa; accordingly,
$c_{\mathrm{head},i}$ should be treated as an experimentally or
EOS-defined leading-arrival speed, rather than assumed to be a single clean
equilibrium longitudinal-wave speed \cite{Hawreliak2024}.

With a selected $c_{\mathrm{head},i}$, calculate
\begin{align}
  \rho_{H,i} &= \rhoz[i]\frac{\Us[i]}{\Us[i]-\up[i]},
      \label{eq:implementation-density}\\
  \CL[{\mathrm{head},i}] &=
      \frac{\rho_{H,i}}{\rhoz[i]}c_{\mathrm{head},i},
      \label{eq:implementation-release-speed}\\
  Z_i &= \rho_{H,i}c_{\mathrm{head},i}
      =\rhoz[i]\CL[{\mathrm{head},i}].
  \label{eq:consistent-implementation}
\end{align}
Use $\CL[{\mathrm{head},i}]$, not $c_{\mathrm{head},i}$, in an
initial-thickness transit estimate.  Use the incremental impedance appropriate
to the actual local release state in Eq.~\eqref{eq:reflection-coefficient}.

Equivalently, a program may store $\CL[i]$ directly.  In that case it must use
$Z_i=\rhoz[i]\CL[i]$, rather than $\rho_{H,i}\CL[i]$.  Either convention is
mathematically consistent; mixing a Lagrangian transit formula with an
Eulerian impedance formula is not.

The following compact sequence is suitable for an overlay:
\begin{enumerate}
  \item Convert all thicknesses to metres and all speeds to SI units.
  \item Infer $u$ from the exact release relation
        Eq.~\eqref{eq:free-surface-exact}, or use $u=U_{fs}/2$ only as a
        declared approximation.
  \item Calculate $x$, $V$, $\Us[f]$, and $\Us[t]$ from
        Eqs.~\eqref{eq:inverse-velocity} and \eqref{eq:linear-hugoniot}.
  \item Obtain a release model $c_i(\sig)$ for quantitative work.  For a
        release-onset estimate, select $c_{\mathrm{head},i}$ and form
        $\CL[{\mathrm{head},i}]$ and $Z_i$ using
        Eqs.~\eqref{eq:implementation-density}--\eqref{eq:consistent-implementation}.
  \item Calculate the free-flight flyer-release time $\tFR$ from
        Eq.~\eqref{eq:flyer-release-free-flight} and the target round trip
        $\Pt$ from Eq.~\eqref{eq:target-round-trip}.
  \item If $\tFR\leq0$, report shock--release catch-up; do not draw a
        sustained plateau.  If $0<\tFR<\Pt$, report release-fan interaction
        and do not draw a simple recompression series.  If $\tFR>\Pt$, report
        that the isolated initial-free-surface return precedes the flyer
        release; $\Pt$ is then an isolated-path diagnostic, not a validated
        fit to the full trace.
\end{enumerate}

All times produced in seconds can be converted to nanoseconds by multiplying
by $10^9$.

\section{First-hand reverberation markers with incomplete target data}
\label{sec:first-hand-markers}

This section defines a deliberately lower-fidelity overlay for experiment
planning and first inspection of a trace.  It must be reported as a
first-hand marker, not as a finite-amplitude prediction.  The indispensable
pre-shot datum is the impact speed $V$.  A peak measured on the subsequently
overlaid trace must not be used to set this prediction.  Flyer identity and
two thicknesses alone do not select a shocked state.

\subsection{Ballpark target $C_0,S$: a state-aware proxy}

Suppose the flyer parameters $(\rhoz[f],C_{0f},S_f)$ are known, while the
target parameters $(\rhoz[t],C_{0t},S_t)$ are only ballpark values for the
identified target material.  With a known $V$, solve the provisional impedance
match
\begin{align}
  \rhoz[t](C_{0t}+S_t\widehat u)\widehat u
  &=\rhoz[f]\bigl[C_{0f}+S_f(V-\widehat u)\bigr]
    (V-\widehat u),
  \label{eq:first-hand-impedance-match}\\
  \widehat x&=V-\widehat u \notag
\end{align}
For positive densities, intercepts, and slopes, the left side of
Eq.~\eqref{eq:first-hand-impedance-match} is strictly increasing in
$\widehat u$ on $[0,V]$, whereas its right side is strictly decreasing
because $\widehat x=V-\widehat u$.  The left side is zero at
$\widehat u=0$ and the right side is zero at $\widehat u=V$; hence there is
exactly one physical root $0<\widehat u<V$.  This supplies the
provisional shock speeds
\begin{equation}
  \widehat{\Us[t]}=C_{0t}+S_t\widehat u,
  \qquad
  \widehat{\Us[f]}=C_{0f}+S_f\widehat x.
  \label{eq:first-hand-shock-speeds}
\end{equation}
The tangent of the assumed target Hugoniot in the stress--particle-velocity
plane gives the Lagrangian \emph{Hugoniot-tangent proxy}
\begin{equation}
  \widehat C_{\mathrm{H},t}^{\mathrm L}=C_{0t}+2S_t\widehat u,
  \qquad
  \widehat C_{\mathrm{H},f}^{\mathrm L}=C_{0f}+2S_f\widehat x.
  \label{eq:hugoniot-tangent-proxy}
\end{equation}
It is neither a shock speed nor, in general, the exact release-head speed.
It follows from $\rhoz^{-1}d\sig_H/d\up$, whereas the true release speed is
given by Eq.~\eqref{eq:release-characteristic-speed}.  Its only role here is
to make a state-aware timing proxy from the available $C_0,S$ values; the
quality of that approximation cannot be inferred from a Hugoniot alone
\cite{LyzengaAhrens1978,Kerley2008}.

When the linear Hugoniot is used in the strength-free sense, this proxy is
most naturally interpreted as a bulk-family surrogate.  For a
strength-bearing solid, a longitudinal release head is generally faster than
the corresponding bulk wave [Eq.~\eqref{eq:longitudinal-bulk-shear}].
Under that additional interpretation, using
$\widehat C_{\mathrm H,t}^{\mathrm L}$ is expected to make
$\widehat P_t$ late relative to a longitudinal-head marker.  This is an
expected one-sided bias, not a rigorous bound: the difference between the
Hugoniot tangent and the actual release adiabat has no universal sign without
a release EOS.  The stated direction is specifically relative to a
\emph{shocked-state} longitudinal head.  It is not preserved if that
comparison is instead made with an ambient longitudinal-speed surrogate:
the tangent proxy can then lie on either side of the surrogate as compression
changes.  A pre-shot longitudinal-speed scenario should therefore be plotted
alongside the tangent-proxy marker whenever suitable shocked-state strength
data are available.

The corresponding first-hand target round-trip and flyer-release markers are
\begin{align}
  \widehat P_t&=\frac{2h_{t0}}{\widehat C_{\mathrm{H},t}^{\mathrm L}}, \\
  \widehat t_{\mathrm{FR}}^{(0)}
  &=\frac{h_{f0}}{\widehat{\Us[f]}}
    +\frac{h_{f0}}{\widehat C_{\mathrm{H},f}^{\mathrm L}}
    +h_{t0}\left(
      \frac{1}{\widehat C_{\mathrm{H},t}^{\mathrm L}}
      -\frac{1}{\widehat{\Us[t]}}\right), \\
  \widehat t_{\mathrm{first}}^{(0)}
  &=\min\bigl(\widehat P_t,\widehat t_{\mathrm{FR}}^{(0)}\bigr).
  \label{eq:first-hand-markers}
\end{align}
Plot $\widehat P_t$ as the candidate onset of an isolated target
free-surface reverberation only when
$\widehat t_{\mathrm{FR}}^{(0)}>\widehat P_t$.  If
$0<\widehat t_{\mathrm{FR}}^{(0)}<\widehat P_t$, flag interacting releases
instead; if $\widehat t_{\mathrm{FR}}^{(0)}\leq0$, flag shock--release
catch-up.  The quantity to compare with the first feature after the initial
free-surface shock state is $\widehat t_{\mathrm{first}}^{(0)}$, not an
unqualified plateau duration.

\subsection{No target $C_0,S$: what can and cannot be estimated}

The flyer Hugoniot cannot provide the missing target release speed.  If target
$C_0,S$ are unavailable, a material name, flyer data, and thicknesses do not
produce a defensible point estimate of $P_t$.  At least one independent target
pre-shot target speed estimate or range must be supplied from a material
library, prior experiment, or literature source; no waveform quantity from
the experiment being overlaid is used in this construction.  If only a
plausible pre-shot Lagrangian timing-speed range
$C_{\mathrm{low}}^{\mathrm L}\leq C_{\mathrm{est},t}^{\mathrm L}
\leq C_{\mathrm{high}}^{\mathrm L}$ is available, report a timing band,
not a fitted point:
\begin{equation}
  \frac{2h_{t0}}{C_{\mathrm{high}}^{\mathrm L}}
  \leq P_t\leq
  \frac{2h_{t0}}{C_{\mathrm{low}}^{\mathrm L}}.
  \label{eq:first-hand-timing-band}
\end{equation}
An ambient bulk or longitudinal value may be used to construct a scenario
range, but it must be labelled as an ambient-speed surrogate rather than a
shocked-state release speed.  A pre-shot absolute shock-arrival estimate also
requires a nominated target shock-speed range,
\begin{equation}
  U_{\mathrm{s},t}^{\mathrm{low}}\leq U_{\mathrm{s},t}
  \leq U_{\mathrm{s},t}^{\mathrm{high}},
  \qquad
  \frac{h_{t0}}{U_{\mathrm{s},t}^{\mathrm{high}}}
  \leq T_s\leq
  \frac{h_{t0}}{U_{\mathrm{s},t}^{\mathrm{low}}}.
  \label{eq:first-hand-shock-arrival-band}
\end{equation}
Without either ballpark target $C_0,S$ or independent pre-shot target speed
ranges, only a qualitative thickness scaling can be reported---not a
quantitative expected trace or reverberation time.

\section{Pre-shot construction of the first-hand expected trace}
\label{sec:pre-shot-overlay}

This procedure uses no measured quantity from the experiment subsequently
overlaid.  Its purpose is to place an expected initial shock state and timing
markers on a plot before the shot; it does not predict the detailed release
shape, pullback amplitude, or fracture state.

\subsection{Inputs and timing markers}

Before the experiment, assemble the flyer density and Hugoniot parameters,
impact speed $V$, flyer and target initial thicknesses, and a target-material
entry.  The target entry should contain either ballpark
$(\rhoz[t],C_{0t},S_t)$ or explicitly sourced pre-shot ranges for target shock
and Lagrangian timing speeds.  The two alternatives are not interchangeable:
the first permits a provisional impedance match, whereas the second produces
only timing bands.

With ballpark target $C_0,S$, use
Eqs.~\eqref{eq:first-hand-impedance-match}--\eqref{eq:first-hand-markers}:
\begin{enumerate}
  \item solve for $\widehat u$ and $\widehat x$ at the impact interface;
  \item calculate $\widehat{\Us[t]}$ and place the initial target
        free-surface shock arrival at
        \begin{equation}
          \widehat T_s=\frac{h_{t0}}{\widehat{\Us[t]}} ;
          \label{eq:pre-shot-shock-arrival}
        \end{equation}
  \item plot the nominal initial free-surface level
        $U_{fs}^{\ast}\simeq2\widehat u$, explicitly labelled as the
        doubling approximation rather than a release-EOS prediction;
  \item place the earliest post-shock timing marker at
        $\widehat T_s+\widehat t_{\mathrm{first}}^{(0)}$;
        draw the isolated target reverberation marker at
        $\widehat T_s+\widehat P_t$ only if
        $\widehat t_{\mathrm{FR}}^{(0)}>\widehat P_t$; and
  \item if $\widehat t_{\mathrm{FR}}^{(0)}\leq\widehat P_t$, annotate the
        plot as an interacting-release or shock--release-catch-up case,
        rather than drawing an $n\widehat P_t$ recompression series.
\end{enumerate}
The displayed plateau is therefore a timing guide only.  Its height is
controlled by the release integral in Eq.~\eqref{eq:free-surface-exact}, and
its later shape requires the finite-amplitude calculation of
Section~\ref{sec:quantitative-reference}.

If target $C_0,S$ are unavailable, do not infer them from the flyer or update
the calculation from the newly acquired trace.  Instead plot the pre-shot
arrival and reverberation bands from
Eqs.~\eqref{eq:first-hand-timing-band} and
\eqref{eq:first-hand-shock-arrival-band}.  If no target shock-speed range is
available, the time origin of the predicted free-surface shock cannot be
placed; if no target Lagrangian-speed range is available, the reverberation
marker cannot be placed.  In either case the appropriate output is a stated
data requirement, not an invented point prediction.

\subsection{Candidate tensile and spall timing window}
\label{sec:spall-window}

Spall is not created merely because a target round-trip time has elapsed.  A
tensile region first requires the interaction of opposed unloading waves, and
damage occurs only if that tensile history exceeds the material's dynamic
spall criterion for a sufficient time.  The interaction of the leading and
trailing release characteristics can nevertheless provide pre-shot
\emph{candidate} times for the tensile window.

Let $T_{I,k}$ be the time at which characteristic $k$ of the flyer-generated
release enters the target at the interface $X=0$, and let $T_{F,j}$ be the
launch time of characteristic $j$ of the target free-surface release at
$X=h_{t0}$.  Let $C_{I,k}^{\mathrm L}$ and $C_{F,j}^{\mathrm L}$ be their
respective Lagrangian speeds through the target.  In a constant-speed
characteristic construction their trajectories are
\begin{equation}
  X_{I,k}=C_{I,k}^{\mathrm L}(T-T_{I,k}),\qquad
  X_{F,j}=h_{t0}-C_{F,j}^{\mathrm L}(T-T_{F,j}).
  \label{eq:opposed-release-rays}
\end{equation}
Their intersection is
\begin{align}
  T_{\cap,kj}
  &=\frac{h_{t0}+C_{I,k}^{\mathrm L}T_{I,k}
      +C_{F,j}^{\mathrm L}T_{F,j}}
      {C_{I,k}^{\mathrm L}+C_{F,j}^{\mathrm L}}, \\
  X_{\cap,kj}
  &=\frac{C_{I,k}^{\mathrm L}
      [h_{t0}+C_{F,j}^{\mathrm L}(T_{F,j}-T_{I,k})]}
      {C_{I,k}^{\mathrm L}+C_{F,j}^{\mathrm L}}.
  \label{eq:opposed-release-intersection}
\end{align}
Use an intersection only when $0<X_{\cap,kj}<h_{t0}$ and the relevant
interface path remains closed.  For positive characteristic speeds, that
interior condition already implies that both characteristics have been
launched.  In an ideal sharp-shock/free-surface construction,
$T_{F,\ell}=T_{F,\tau}=T_s$, but a centered release fan already contains
distinct leading ($\ell$) and trailing ($\tau$) characteristics with
different speeds.  A finite shock rise adds a spread in their launch times;
a state-dependent release determines their speed spread.

Let $\mathcal I$ contain the physically admissible leading--trailing
intersections from Eq.~\eqref{eq:opposed-release-intersection}.  The
release-overlap window is
\begin{equation}
  T_{\mathrm{overlap,first}}^{\ast}
  =\min_{(k,j)\in\mathcal I}T_{\cap,kj},
  \qquad
  T_{\mathrm{overlap,last}}^{\ast}
  =\max_{(k,j)\in\mathcal I}T_{\cap,kj}.
  \label{eq:spall-window-markers}
\end{equation}
For the usual monotonic fan ordering, these extrema are the head--head and
trailing--trailing intersections, respectively.  The explicit minimization
and maximization avoid assuming that ordering for an arbitrary release model.
The second marker is the end of a coarse release-fan-overlap window, not a
universal latest spall time.  With an instantaneous fracture criterion, spall
could begin no earlier than $T_{\mathrm{overlap,first}}^{\ast}$ if the tensile
threshold is crossed.  Real fracture can be delayed by incubation, void
nucleation, and growth; hence a defensible last spall-formation time requires
a strength/damage model and cannot be obtained from leading and trailing wave
speeds alone \cite{Hawkins2020}.

For an expected free-surface trace, a pullback created at an estimated
intersection plane must still propagate from that plane to the target free
surface.  A first arrival marker is therefore bounded by
\begin{equation}
  T_{\mathrm{pb,first}}^{\ast}\gtrsim
  T_{\mathrm{overlap,first}}^{\ast}
  +\frac{h_{t0}-X_{\mathrm{overlap,first}}}
         {C_{\mathrm{pb}}^{\mathrm L}},
  \label{eq:pullback-first-arrival}
\end{equation}
where $X_{\mathrm{overlap,first}}$ is the intersection position associated
with $T_{\mathrm{overlap,first}}^{\ast}$ and
$C_{\mathrm{pb}}^{\mathrm L}$ is the appropriate post-intersection signal
speed.  This is distinct from the
intact-target marker $\Pt$.

The single Hugoniot-tangent proxy in
Eq.~\eqref{eq:hugoniot-tangent-proxy} provides only one characteristic speed;
it cannot determine the leading--trailing fan width in
Eq.~\eqref{eq:spall-window-markers}.  A pre-shot spall-window scenario
therefore needs independently chosen leading/trailing release-speed ranges
and release-launch durations from prior material information.  Without those
inputs, report only the nominal release-intersection marker and do not label
it a spall-time prediction.

\section{Limits of the first-hand overlay}

The first-hand construction intentionally omits the information required by
the quantitative reference calculation of Section~\ref{sec:quantitative-reference}.
It should not be used to identify damage or to tune a material speed until a
marker coincides with a chosen feature.  In particular, a quantitative
finite-amplitude plate-impact trace still requires:
\begin{itemize}
  \item a release adiabat or EOS (for example, a calibrated Mie--Gr\"uneisen
        model) for each material, rather than only $C_0$ and $S$;
  \item shocked-state longitudinal sound speeds along that release path;
  \item a strength/elastic--plastic model if precursor or strength effects are
        material to the trace;
  \item characteristic interactions at the flyer--target interface and both
        free surfaces, together with a unilateral contact condition if the
        flyer is unbonded; and
  \item a separate spall model if a pullback/recompression is attributed to a
        spall plane.  A spall-plane recompression is not the same phenomenon as
        the no-spall target reverberation of Eq.~\eqref{eq:initial-fs-reverberation}
        \cite{Hawkins2020}.  For example, if a spall plane is at initial
        material-coordinate depth $\delta$ below the target free surface, its simplest local
        round-trip marker is $2\delta/\CL[t]$, rather than the intact-target
        marker $\Pt=2h_{t0}/\CL[t]$; the actual timing still depends on the
        wave paths and evolving material state.
\end{itemize}

Accordingly, a model based only on $\rhoz$, $C_0$, $S$, an ambient longitudinal
speed, and two thicknesses should be presented as a first-order diagnostic
timing guide, not as an exact prediction of a spall trace
\cite{Kerley2008,Hawreliak2024,Hawkins2020}.

\begin{thebibliography}{9}
\bibitem{Kerley2008}
G. I. Kerley, \emph{Calculation of Release Adiabats and Shock Impedance
Matching}, Kerley Technical Services Report KTS08-1 (2008).
\url{https://arxiv.org/pdf/1306.6913}

\bibitem{Kerley2006}
G. I. Kerley, ``The Linear $U_S$--$u_P$ Relation in Shock-Wave Physics,''
Kerley Technical Services Report KTS06-1 (March 2006; arXiv:1306.6916, 2013).
\url{https://doi.org/10.48550/arXiv.1306.6916}

\bibitem{Rice1958}
M. H. Rice, R. G. McQueen, and J. M. Walsh, ``Compression of solids by strong
shock waves,'' in \emph{Solid State Physics}, Vol.~6, edited by F. Seitz and
D. Turnbull (Academic Press, New York, 1958), pp.~1--63.
\url{https://doi.org/10.1016/S0081-1947(08)60724-9}

\bibitem{LyzengaAhrens1978}
G. A. Lyzenga and T. J. Ahrens, ``The relation between the shock-induced
free-surface velocity and the postshock specific volume of solids,''
\emph{Journal of Applied Physics} \textbf{49}, 201--204 (1978).
\url{https://doi.org/10.1063/1.324323}

\bibitem{Hawreliak2024}
J. A. Hawreliak, J. M. Winey, Y. Toyoda, K. Zimmerman, and Y. M. Gupta,
``Sound speed determination in copper shock compressed to 190 GPa,'' \emph{Journal of
Applied Physics} \textbf{136}, 165902 (2024).
\url{https://doi.org/10.1063/5.0220264}

\bibitem{CourantFriedrichs1948}
R. Courant and K. O. Friedrichs, \emph{Supersonic Flow and Shock Waves}
(Interscience, New York, 1948).

\bibitem{Davison2008}
L. Davison, \emph{Fundamentals of Shock Wave Propagation in Solids}
(Springer, Berlin and Heidelberg, 2008).
\url{https://doi.org/10.1007/978-3-540-74569-3}

\bibitem{Hawkins2020}
M. C. Hawkins, S. A. Thomas, S. J. Fensin, D. R. Jones, and R. S. Hixson,
``Spall and subsequent recompaction of copper under shock loading,''
\emph{Journal of Applied Physics} \textbf{128}, 045902 (2020).
\url{https://doi.org/10.1063/5.0011645}
\end{thebibliography}

\end{document}
